% Source: https://artofproblemsolving.com/community/c3h1433412_geometry_question
% Stable identity: aops-topic-1433412-q01
\textbf{Answer:} $BX:XM=2a^2:b^2$.

\begin{asy}
usepackage("amsmath");
size(5.5cm);
defaultpen(fontsize(10pt));
pair ft(pair P,pair U,pair V){return U+dot(P-U,V-U)/dot(V-U,V-U)*(V-U);}
path ra(pair P,pair u,pair v,real s=0.2){return P+s*unit(u)--P+s*unit(u)+s*unit(v)--P+s*unit(v);}
real a=3,bb=4; pair C=(0,0),A=(bb,0),B=(0,a);
pair D=ft(C,A,B), M=(A+C)/2, X=extension(C,D,B,M), K=(A+D)/2;
draw(M--K,deepgreen+dashed); dot(K); label("$K$",K,NE);
draw(A--B--C--cycle,linewidth(0.9));
draw(C--D,blue); draw(B--M,red);
draw(ra(C,A-C,B-C,0.25)); draw(ra(D,A-D,C-D,0.2),blue);
dot(A);dot(B);dot(C);dot(D);dot(M);dot(X);
label("$A$",A,SE); label("$B$",B,N); label("$C$",C,SW); label("$D$",D,NE); label("$M$",M,S); label("$X$",X,dir(150));
label("$a$",(0,a/2),W); label("$\dfrac{b}{2}$",(C+M)/2,S); label("$\dfrac{b}{2}$",(M+A)/2,S);
label("\mbox{}",(min(currentpicture,user=true)+max(currentpicture,user=true))/2,black); // keeps the text after the figure black when LaTeX typesets the labels inline
\end{asy}

\textbf{Solution with similar triangles.} Let $c=AB$. Triangle $ADC$ is similar to triangle $ACB$ (they share angle $A$ and have right angles at $D$ and $C$), so $\dfrac{AD}{AC}=\dfrac{AC}{AB}$, that is $AD=\dfrac{b^2}{c}$; similarly $BD=\dfrac{a^2}{c}$. Draw the line through $M$ parallel to $CD$, meeting $AB$ at $K$. Since $M$ is the midpoint of $CA$ and $MK\parallel CD$, $K$ is the midpoint of $AD$, so $DK=\dfrac{AD}{2}=\dfrac{b^2}{2c}$. In triangle $BMK$, $XD\parallel MK$, so
\[
\dfrac{BX}{XM}=\dfrac{BD}{DK}=\dfrac{a^2}{c}\cdot\dfrac{2c}{b^2}=\dfrac{2a^2}{b^2}.
\]
\textbf{Check with coordinates.} Put $C=(0,0)$, $A=(b,0)$, $B=(0,a)$. The altitude from $C$ has direction $(a,b)$ (perpendicular to $\overrightarrow{BA}=(b,-a)$), and the median is $B+t(M-B)=\left(\dfrac{tb}{2},\ a(1-t)\right)$. This point lies on the altitude when $\dfrac{tb}{2}\cdot b=a(1-t)\cdot a$, that is $t=\dfrac{2a^2}{2a^2+b^2}$. So $BX:XM=t:(1-t)=2a^2:b^2$.

So the ratio depends on the shape of the triangle: for example it is $2:1$ for an isosceles right triangle, and $18:16=9:8$ for legs $a=3$, $b=4$ (the figure).
